We explain the character and group-order bounds used here.
Over the residue field \(\mathbb F_Q\) of cardinality \(Q\), the mirabolic subgroup is
\(\operatorname{GL}_{r-1}(\mathbb F_Q)\ltimes\mathbb F_Q^{r-1}\).
Counting successive independent columns and using

\[
  |U_Q|=Q^{r(r-1)/2}
\]
 shows that its index divided by
\(|U_Q|\) is \(\prod_{a=1}^{r-1}(1-Q^{-a})\). This formula
also covers the empty product. For the product over the
two level primes the ratio \(I/|U|\) is at least
\(1-\sum_Q 1/(Q-1)\), since the sum of all the powers
\(Q^{-a}\) at one place is at most \(1/(Q-1)\). There
are at most \(2f\) places, and each \(Q\) is at least
\(\min(g_0,g_1)\). The common size record has \(f\le k\) and
\(\min(g_0,g_1)\ge(kb)^8\), so
\(2f/(\min(g_0,g_1)-1)<1/8\).

\begin{lemma}[The character dimension at the definite lift]
\label{lem:paired-character-share}
For the definite two-level cusp model at the selected weight, write \(V\) for the full cusp space of dimension \(R_j\). Its character multiplicity at the given level is
\[
 \mu=(\chi|_U)\psi_0^{-1},\qquad \dim V[\mu]=\frac{R_j}{|U|}.
\]

\end{lemma}

\begin{proof}
\proofstep{Free orbits at the definite lift.}
For the dimension bound use the two-level model of cusp
forms for this same lift. Scalar extension through the
selected complex embedding preserves the ambient dimension
and the dimension of the character space. In fact the
images of the coordinate basis are independent and span,
and the joint eigenspace is the kernel of the stacked
matrices \(\rho(u)-\mu(u)\) for \(u\in U\).
Equivalently, if \(V\) is the model space, the selected subspace is
\[
  V[\mu]=\bigcap_{u\in U}\ker\big(\rho(u)-\mu(u)I\big).
\]
 All their
nonzero minors stay nonzero. Use the model character
\(\mu=(\chi|_U)\psi_0^{-1}\), including the scalar
twist on the model action.

At the definite end the
subgroup acts freely on the components at this level.
A stabilizer would give a unitary rational matrix at a
balanced frame, unipotent modulo the specified split prime
and bounded in the positive definite norms. The
small-element estimate at that level forces it to be the
identity, and faithfulness gives the identity subgroup
element. The direct sum of the component spaces is therefore
a sum of regular representations of \(U\), by identifying
a representative fiber with each of its translates. Its
\(\mu\)-space has dimension \(R_j/|U|\).
\label{cor:paired-character-share-1}

\end{proof}

\subsubsection*{Equivariant realization of the selected row as cusp forms}
\label{subsec:selected-cusp-reading}

An index computed in geometric coordinates controls the selected row only through an equivariant identification. We exhibit the cusp-form realization on which the local cuts and the level actions agree.

Fix a lift $j$ and the run used at a coefficient place $v$.
Write $K_0$ for the construction field of the run, $K$ for
the common row field and $C_j$ for the rational cusp space
at its prescribed weight $\kappa$. The form coordinates
identify $C_j$ with a $K_0$-subspace of $K^{R_j}$ and
identify $K\otimes_{K_0} C_j$ with $K^{R_j}$. If \(b_1,\ldots,b_d\) is the common basis in the row data over \(K_0\), scalar extension preserves every coordinate independently:
\[
  K\otimes_{K_0}\bigoplus_{i=1}^dK_0 b_i=\bigoplus_{i=1}^dK(1\otimes b_i).
\]
 The chosen complex place determines an embedding
$\sigma_0:K_0\hookrightarrow\mathbb C$ and an extension
$\sigma:K\hookrightarrow\mathbb C$. Use the bijection
between the two-level components and the run's components.
At corresponding components the Hermitian data, frame
domains, embeddings and weights agree. Restrict the holomorphic forms to these frame domains.

The ordinary stable-range theta calculation in
Section~\ref{subsec:theta-row-rank} supplies the span before
projection. Here the target has rank $r$ over $E$ of
degree $2f$, the positive definite source has even rank
$n_V\ge 8fr+4$, and the level is odd with chosen split half
coprime to its conjugate.
The degree and rank conditions in the stable range combine to give
\[
  n_V\ge 8fr+4=4[E:\mathbb Q]r+4 .
\]
 Use the run's prescribed first
character Gram comparison and its admissible weight,
including positivity of the last weight. The chart
independence and the short theta transformation specialize
the stable unfolding to readings of raw short theta
sections. Every such short section is cuspidal at this
weight.

On each component extend a bounded holomorphic
section by zero off its frame domain, as a function on
all frame coordinates. The extension belongs to the
space on which the theta span is taken. The stable unfolding
expresses its family of components in a finite complex
span of the read short sections. Restricting the resulting
identity to the domains and transporting by the component
bijection gives the desired expansion of the original
family on the two-level model. The construction uses
the positive bad modulus and the Gram comparison of the
actual run.\label{cor:paired-row-reading-1}

\begin{proposition}[The equivariant cusp-form realization and the rank at the smaller level]
\label{prop:equivariant-selected-row}
For a lift \(j\) of either run in
Proposition~\ref{prop:linked-comparison-rows}, use its rational
cusp space \(C_j\) at weight \(\kappa\) over \(K_0\), with the
form coordinates of Definition~\ref{def:common-row-lattices}.
Extend along the complex embedding \(\sigma_0\) of \(K_0\). The chosen basis extends to a basis of the complex space of holomorphic cusp forms and identifies the corresponding character actions. The retained arithmetic projector is the same
one in the two realizations. For each lift the partition of the retained base labels gives the small-level image dimension
\[
 dZ_0(j)=D_0.
\]
 The full rank of the integral image uses the separate
off-prime trace inclusion.
\end{proposition}

\begin{proof}
\proofstep{Independence in the two coefficient realizations.}
For independence take a finite $K_0$-independent family of
cusp rules. Evaluate it on the framed test algebras through
all maps to $\overline{K_0}$. These evaluations detect a
relation of rules. Indeed the test algebras are reduced
and of finite type, so a nonzero test value survives at
a closed point over a finite extension of $K_0$. The span
of the evaluation rows is invariant under
$\operatorname{Gal}(\overline{K_0}/K_0)$, because composing
each test map with an automorphism retains it in the family.

If the row rank were deficient, the unique reduced row
echelon basis of the annihilator would be Galois invariant.
A nonzero basis row would give a $K_0$-relation among the rules. Hence some set \(I\) of rows, as large as the family of rules, gives a square matrix \(A_I\) with nonzero determinant. Under the chosen complex embedding \(\sigma\),
\[
  \det\big(\sigma(A_I)\big)=\sigma(\det A_I)\ne0 .
\]
 Framed base change realizes these tests
as complex frames. Uniformization and the common Borel
weight multiply all entries of an evaluation row by the
same nonzero scalar when passing to its frame-domain
representative. The resulting functions on frames are therefore independent. Restriction to the model domains preserves independence, since the functions vanish outside these domains and every component occurs.\label{cor:paired-row-reading-2}

\proofstep{The same rational character space on the two levels.}
Read the common $K_0$-basis in the two-level model and
extend its coefficients by $\sigma$. Independence and the
stable span give the isomorphism
\[
 \mathbb C\otimes_{K,\sigma} K^{R_j}\ \simeq\ V_\kappa .
\]
On a rational cusp form the lift of a level change acts on
the right in its moduli coordinate. In frame coordinates
this is the two-level pullback, with the same scalar and
the same change of component as in the run. The equality
on the basis therefore intertwines the full level action
after extension of scalars. In particular the dimensions
of the character kernel and its intersection with the kernel
of the good projector can be computed in $V_\kappa$.
Let \(A\) be the stacked character equations and \(P\) the good projector in these coordinates. Their common kernel is computed by the single matrix
\[
  \ker A\cap\ker P=\ker\begin{pmatrix}A\\P\end{pmatrix}.
\]
Changing the complex embedding preserves their algebraic
ranks.\label{cor:paired-row-reading-3}

\proofstep{Comparing with the retained arithmetic projector.}
The identification with the retained arithmetic rank uses
one more comparison at the selected point. In its row
coordinates the tie is the composite of the map to
$K\otimes_{K_0}C_j$ with the form-coordinate isomorphism
above. It identifies the subgroup of the arithmetic row
with the preimage of the base unipotent subgroup under the
level injection and identifies their characters there.
The determinant character in the coordinate change is
one on that subgroup. Consequently the arithmetic character
space maps bijectively to $V_\lambda$ after extension of
scalars.
Consequently the complex identification preserves the character rank, since scalar extension commutes with the common kernel
\[
  \bigcap_{u}\ker\big(\rho(u)-\mu(u)I\big).
\]
 The second row link uses the base-tuple theta
lift and its rational reading. In the split part of the
chosen realization these give the same reduced level
group and the same target projector, because the scalar
phase at the base is one.

The base-level lift identifies
the retained source space with
$V_\lambda\cap\operatorname{im}e_{\rm good}$ in the target
row. Its source dimension is the common $D_0$ of the
point, for every lift $j$. In the rank comparison the
base-tuple maps identify the dimensions of the summands
indexed by each retained irreducible type. Those summands
exhaust the respective spaces, so the sums of their
dimensions agree. Thus the target character rank
is the arithmetic $d_0$, and the retained rank is
\[
 D_0=d_0-\dim(V_\lambda\cap\ker e_{\rm good}).
\]
The equality uses invariance of the character space under
$e_{\rm good}$ and its decomposition into the image and
kernel of the restricted idempotent. In the integral
comparison all row lattices are transported by
this same tie. The scalar extension here computes ranks
and the lattice comparison uses localization in the
number field.\label{cor:paired-row-reading-4}

\proofstep{The rank in the determinant denominator.}
The small-level rank in the division denominator can be
read from the projection at the base tuple. The level
projection at the point's base index \(t_0\) is surjective
for each lift. The tuple readout sends \(t_0\) to \(\mu_0\),
since the base Hecke character is the identity and the
tuple characters distinguish the used indices.
On the matched row models the level projection induces the
theta adjoint \(\Theta_{\mu_0,j}^{*}\) by the class-sum identity
and the two level actions. Thus this adjoint surjects onto
the reference row space. Since

\[
  \ker\Theta_{\mu_0,j}=(\operatorname{im}\Theta_{\mu_0,j}^*)^\perp,
\]
\(\Theta_{\mu_0,j}\) is injective.

Decompose the reference row into the occurring class types.
The tied eigenspace multiplicity formula for each type gives
the sum of the row counts over that type's record labels.
These label sets partition the \(D_0\) base row labels,
and the class type summands form an internal direct sum.
Summing gives \(D_0\) at every lift. By injectivity the
small-level range \(\operatorname{im}\Theta_{\mu_0,j}\)
has that same dimension. This proves the identity
\(dZ_0(j)=D_0\) used in the cohomological division term.
The full rank of the projected \(J_v\) at a twisted index
follows separately from the local trace inclusion on the
assigned tuple.
\label{cor:paired-row-reading-5}
\end{proof}

\paragraph{The selected levels and the separating section.}

Equivariance identifies the row as a cusp space. Separating its level components also requires a section that distinguishes those components at the chosen embedding.

The components are those of the complex uniformization of the chosen construction at its fixed complex place. A component $c$ has lower
arithmetic group $\Gamma_c$ and its two prescribed lattices.
The change of level sends $c$ to $hc$ and has a hermitian
transporter $t(h,c)$ carrying both lattices to their counterparts.
The transporter multiplication is well defined modulo the lower group of the destination. For \(h,k\) in the level group it gives
\[
  t(h,kc)t(k,c)t(hk,c)^{-1}\in\Gamma_{hkc}.
\]
 A transporter fixing its
component lies in the lower group only for the identity level element.
The lower character on the arithmetic group is trivial. The
components have the common signature, CM embeddings and weight.
Use the equivariant cusp-form realization of
Section~\ref{subsec:selected-cusp-reading} at this weight. It identifies
the scalar extension of the whole row with the space of bounded
holomorphic sections on these components.
\label{cor:paired-cusp-ranks-1}

Let \(H\) reduce to upper unitriangulars at one chosen level prime \(g\), and let \(H_c\) stabilize the component \(c\). Take the upper group \(G\) generated by \(\Gamma_c\) and the transporters of \(H_c\). The transporter product law above and the trivial kernel on level elements give
\[
  1\longrightarrow\Gamma_c\longrightarrow G\longrightarrow H_c\longrightarrow1.
\]
 It preserves both lattices and is unipotent modulo every prime
of $E$ above $g$. The congruence on the chosen half of the primes
follows from the level change and the kernel congruence for $\Gamma_c$.
On the conjugate half, the reduction of a unitary matrix is the
contragredient with respect to the perfect paired form, so it too
is unipotent.
If \(A\) is unipotent, so is its inverse transpose on the conjugate half, since
\[
  A^{-\mathsf t}-I=-A^{-\mathsf t}(A^{\mathsf t}-I),\qquad (A^{\mathsf t}-I)^r=0 .
\]
 The prime $g$ splits completely in $E$ and exceeds $3r$.
These conditions apply to the character subgroup, the two half-level
subgroups used below, and the trivial subgroup.
\label{cor:paired-cusp-ranks-2}

\begin{lemma}[Separation at the actual two-level translates]
\label{lem:separating-level-section}
Let \(x_b\) denote the coset centers in the chosen cut and \(c\) the selected center. There is a cusp section \(s_c\) of the prescribed determinant weight whose values in the chosen frames satisfy
\[
 s_c(x_b)=\begin{cases}1,&b=c,\\0,&b\ne c.\end{cases}
\]
 The comparison uses the marked subgroups at the two
actual levels and the same coefficient realization as the
character calculation.
\end{lemma}

\begin{proof}
\proofstep{A cusp section distinguishing the level translates.}
Use a balanced majorant of the self-dual component lattice, in a
frame $x_c$ for the Hermitian form and that majorant. The property
needed of the majorant is that an endomorphism of the lattice
strictly contracting at every complex embedding is zero. The
semistable-majorant statement gives this property by applying the
degree inequality to the kernel and image of a contraction. The lattice has degree zero at the semistable majorant, and its
sublattices have degree at most zero. A nonzero strict contraction
would increase the degree of the quotient by its saturated kernel
when mapped isomorphically onto its image lattice. This is
incompatible with the degree bound for that image in the lattice.
Indeed for a nonzero strict contraction \(T\) the saturated kernel gives the incompatible inequalities
\[
  0\le\widehat{\deg}(L/\ker T)<\widehat{\deg}(T L)\le0 .
\]
 The property is invariant under the
right triangular change to an orthonormal frame.

Here is the small-element estimate in this frame. If the operator
norm of $\gamma$ is at most $e^{1/(3r)}$ at every embedding, then
for $1\le a\le r$,
\[
 \operatorname{tr}(\gamma^a)-r\in g\mathcal O_E,
 \qquad |\tau(\operatorname{tr}(\gamma^a)-r)|
       \le r(e^{1/3}+1)<g .
\]
The trace congruence uses nilpotence of the reduction of $\gamma-1$
at each prime over $g$ and squarefreeness of $g\mathcal O_E$. The
field norm now forces all the displayed trace differences to vanish.
Newton's identities give the characteristic polynomial $(X-1)^r$.

In an orthonormal basis the squared Hilbert--Schmidt norm of
$\gamma-1$ is at most $r(e^{2/(3r)}-1)\le e^{2/3}-1<1$, using
$\operatorname{tr}\gamma=r$.
The Hilbert--Schmidt computation uses the trace to cancel both linear terms, giving
\[
  \|\gamma-I\|_{\rm HS}^2=\operatorname{tr}(\gamma^*\gamma)-r\le r(e^{2/(3r)}-1)<1 .
\]
 Its operator norm is strictly below
one, so balance gives $\gamma=1$. The radius potential $R$ on
the positive-plane domain measures the change of majorant from the
orthonormal center. If two points of the ball $R<1/(49r^2)$ are
related by $\gamma$, their two Gram comparisons bound the norm
of $\gamma$ at the center by $e^{1/(3r)}$. This gives a ball whose
nonidentity upper arithmetic translates are disjoint. In the
definite case all upper arithmetic matrices are isometries in the
majorant, so the upper group itself is trivial.
\label{cor:paired-cusp-ranks-3}

For $0<p_0<r$ put $n=f r(r-1)/2$, $d_X=f p_0(r-p_0)$ and
\[
  s_0=2r+50r^2d_X+2,\qquad d=\tfrac12s_0(-1^{[p_0]},+1^{[r-p_0]}).
\]
Here the two signed blocks are repeated at all CM embeddings.
We use the separating section of weight $d$ at the free center.
For clarity, the relevant orbit interpolation estimate concerns an
upper arithmetic group $G$ containing $\Gamma_c$, a free ball
$R<\rho_0=1/(49r^2)$, and the weight potential $\varphi_d$.
The signed two-block calculation used in \hyperref[par:slice-interpolation]{the cut interpolation}
gives $\mathcal L(\varphi_d-\log D)\ge (s_0-2r+2)\mathcal L(R)$.
This margin pays for a radial weight $bR$ and for logarithmic
poles of order $d_X$ at the orbit, since
\[
  s_0-2r+2 > b+d_X/\rho_0.
\]

When \(p_0(r-p_0)\ge4\), use the exponent
\[
  b=4fr^2+1.
\]
 For the smaller signatures
one uses the sum estimate on the positive-plane domain with
$b=4(r-1)$, since the compact flag fibers contribute no factor in
this block-constant weight. In either range the corresponding choice of \(b\) gives a constant \(C\) for the individual model with
\[
  \sup_x\sum_{\gamma\in\Gamma_c}\exp\!\left(-\tfrac b2 R(\gamma x)\right)\le C.
\]
 These bounds use separation in
the orthogonal coordinates and integration of the invariant
volume density over the disjoint balls. All these constants are
for the individual model and weight. Only the stated weight shift
will be used across the family.
\label{cor:paired-cusp-ranks-4}

\proofstep{Interpolation with simultaneous vanishing on the other rows.}
On the flag chart take a smooth cutoff equal to one near the center
and supported inside the free ball. Place smoothed logarithmic
poles at the $G$-orbit of the underlying positive plane. The poles
are truncated to zero outside the disjoint balls. Their Levi loss
is at most $(1+\eta)d_X/\rho_0$ times $\mathcal L(R)$ for a fixed
sufficiently small $\eta>0$.
For the smoothed pole potential \(P\), the curvature cost on these balls is
\[
  {\cal L}(P)\ \ge\ -\frac{(1+\eta)d_X}{\rho_0}\,{\cal L}(R).
\]
 The smoothing parameter tends to zero, with the potential
uniformly bounded on the cutoff's derivative annulus.
Solve $\bar\partial$ of the cutoff with total weight
$\varphi_d-\log D-bR$ plus this logarithmic potential. The same
weighted estimate and frame extension as for the cut interpolation
apply. The solution has uniformly bounded weighted norm as the
smoothings tend to their poles.

In a small ball in the $d_X$ transverse
directions around any orbit point, the weight in the squared-norm integral
is bounded below by a fixed multiple of
$(|z|^2+\varepsilon)^{-d_X}$ times the inverse of the radial
factor. The radial mass diverges logarithmically. For \(d_X>0\) and
a fixed small \(\delta_0>0\), it satisfies
\[
 \int_0^{\delta_0}\frac{r^{2d_X-1}}{(r^2+\varepsilon)^{d_X}}\,dr
    \ \ge\ 2^{-d_X}\int_{\sqrt\varepsilon}^{\delta_0}\frac{dr}{r}
    \ \longrightarrow\ \infty .
\]
Near each orbit centre the cutoff is locally constant, so the
correction there is holomorphic. Its squared sphere mean bounds
its squared value at the centre. The uniform weighted norm bound,
divided by this diverging mass and with the radial weight evaluated
at that centre up to the common comparison constant, makes the
correction at the centre tend to zero. The
bound at each other $G$-orbit point has the additional factor
$e^{-bR/2}$ coming from the radial weight at that point.

Thus after
Poincar\'e summation, the cusp section has readings arbitrarily close
to one at the center and zero at all other coset centers.
For a sequence of such sums \(F_n\), with the representatives indexed by \(b\), the limiting readings have the form
\[
  \lim_{n\to\infty}F_n(x_b)=\begin{cases}1&b=c,\\0&b\ne c.\end{cases}
\]
 Summation
uses the orbit bounds just stated. The bounded section is cuspidal
in this block-constant weight with positive cross-block gap, by the
cusp translation and decay criterion. The map from cusp sections
to the finite tuple of readings has a closed image, being a linear
subspace of a finite-dimensional space. It contains the exact delta
tuple.
Writing \(E\) for the finite tuple evaluation and \(\delta_c\) for the desired delta tuple, closedness yields
\[
  \delta_c\in\overline{\operatorname{im}E}=\operatorname{im}E .
\]
 This gives the separating section on a component, extended by
zero on the other components. At a definite component the upper
group is trivial and the determinant power itself is nonzero there.
\cite{hormander-slice}\label{cor:paired-cusp-ranks-5}

\proofstep{Returning the separating section to the chosen components.}
For each $H$-orbit of components choose a representative $c$ and
its separating section. Sum these sections over the orbits. For the
sum $u_0$, the reading of $g^*u_0$ at $x_c$ vanishes unless $g=1$.
Inside the stabilizer this follows from the separation modulo
$\Gamma_c$, and outside it the component supports are disjoint.
At the transported frame of $kc$ the entry $(h,g)$ of the translate
matrix reads $(hgk)^*u_0$ at $x_c$. Solving the group equation gives the unique nonzero entry of row \(h\):
\[
  hgk=1\quad\Longleftrightarrow\quad g=h^{-1}k^{-1}.
\]
 The determinant therefore has a
nonzero value on every component, including when the group is
noncommutative. This construction is available independently for
all the indicated subgroups with the same weight $d$.
\label{cor:paired-cusp-ranks-6}

\end{proof}

\begin{established}[The flat cut and the determinant-weight gaps]
\label{hyp:three-determinant-weight-ranks}\label{cor:paired-cusp-ranks-7}
At an indefinite lift write \(n=f\binom r2\), \(T_0=T\),
\(L=w_L\), \(F=w_F\), and \(\kappa=F+T_0 L\). On the common
flat cut of \hyperref[par:slice-interpolation]{the cut interpolation},
\(n\ge2\), \(11n^2\le T_0\), and \(s_0=2r+50r^2d_X+2\le T_0\).
Number the flag coordinates \(0,\ldots,r-1\), negative block first,
and put \(k_{\rm vol}(i)=2i-(r-1)\). For
\[
                   U=F-k_{\rm vol}-(n+2)L
\]
the adjacent decreasing gaps inside each block are at least zero,
and \(U_b-U_a\ge1\) for every negative coordinate \(a\) and
positive coordinate \(b\). For \(L\) the corresponding lower
bounds are \(1\) and \(19T_0/10\), respectively.
\end{established}
\begin{proof}
The same light columns and auxiliary section give the flat
intersection for the three weights by
\hyperref[par:slice-interpolation]{the cut interpolation}.
Lemma~\ref{hyp:common-slice-scales} gives \(T_0\ge11n^2\);
\(d_X=fp_0(r-p_0)\le n\), and the polynomial bounds in its proof
give \(2r+50r^2n+2\le T_0\) at large scales.

Set \(d=n+2,\ T'=T_0+d\). For an internal gap \(a\) of the
unsigned row of \(\kappa\), the light gap \(l\) is defined by
\(a+2=T'l+\rho,\ 0\le\rho<T'\).
Lemma~\ref{hyp:height-section-weights} gives \(l\ge1\).
Since \(F=\kappa-T_0L\), after reversing and negating the negative
block the increasing shift in the volume weight contributes two:
\[
 a+2=T'l+\rho,\qquad k_{\rm vol}(i+1)-k_{\rm vol}(i)=2,\qquad
 U_i-U_{i+1}=a-(T_0+d)l+2=\rho\ge0 .
\]

At an unsigned terminal entry write
\(a-d_0=T'l+\rho\), where the two offsets \(d_0\) are \(dn_W+2r\)
and zero. Lemma~\ref{hyp:height-section-weights} gives \(a\ge2T_0^2\)
and \(T'\le41T_0/40\). Consequently
\[
 l>\frac{2T_0^2-T_0}{41T_0/40}-1\ \ge\ \frac{19}{10}T_0
 \quad(T_0\ge 100).
\]
The light rows have nonnegative central shifts of opposite signs,
so every positive coordinate minus every negative coordinate has
at least this gap. The unsigned terminal entries of \(F-dL\) are
\(d_0+\rho\). In the difference of the last positive and the first
negative flag coordinate the central shifts of \(F-dL\) contribute
\(n_P-dn_W\); subtracting \(k_{\rm vol}\) costs \(2(r-1)\).
The resulting difference for \(U\) is at least
\[
 (dn_W+2r)+0+n_P-dn_W-2(r-1)=n_P+2\ge1 .
\]
Internal monotonicity makes this the smallest cross-block difference.
\end{proof}

\begin{proposition}[Dimensions at the three determinant weights]
\label{prop:three-selected-weight-ranks}
Use \(n\ge2\), \(T_0\ge11n^2\), and one common flat cut
of Lemma~\ref{hyp:three-determinant-weight-ranks}, with
\(\delta_c\) points on component \(c\). Put
\(\delta=\sum_c\delta_c\). Then the cusp dimensions satisfy
\[
 \dim V_{\kappa-d}\ge\delta\binom{T_0+n}{n},\qquad
 \dim V_{\kappa}\le\delta\binom{T_0+2n+6}{n},\qquad
 \dim V_{\kappa+2d}\le\delta\binom{T_0+3n+10}{n}.
\]
In particular \(\dim V_{\kappa+2d}\le 2\dim V_{\kappa-d}\). The dimensions used in the comparison
are those of the prescribed cusp spaces.
\end{proposition}

\begin{proof}
\proofstep{The lower rank by interpolation.}
For the lower count, subtract $d$ from $F$. In the interpolation
weight $F-d-k_{\rm vol}-nL$ the cross-block coefficient is at least
$1+3.8T_0-s_0$. This exceeds $4fr^2$, since $fr^2\le4n^2$.
Subtracting one further $L$ leaves the cross-block coefficient
positive, with all within-block coefficients nonnegative. These
are the curvature and additional margin required by the cut
interpolation with zero additional potential. It gives bounded
sections $Q_j$ of weight $F-d$ whose readings at the cut are the
coordinate unit vectors.

For each reading $j$ and $|\alpha|\le T_0$,
form
\[
   Q_j\,w_0^{T_0-|\alpha|}\prod_{a=1}^n S_a^{\alpha_a}.
\]
At least one light factor is cuspidal, including when $|\alpha|=T_0$.
In a linear relation, take the least surviving total degree at the
points and use the local coordinates $S_a/w_0$. Their leading
coefficients read the corresponding $Q_j$ alone. They vanish only
when all the coefficients at that degree vanish, so the family is
independent. For $\delta_c$ cut points the lower dimension at
$\kappa-d$ is $\delta_c\binom{T_0+n}{n}$.
\label{cor:paired-cusp-ranks-8}

\proofstep{The upper rank by injecting into jets.}
For the upper count, the jet map at these same points is injective
through degree $T_0+n+6$. For \(f_0\) of weight \(F+vL\), put
\[
 Q=\frac{\|f_0\|^2\|\theta_2\|^2\|\theta_1\|^2\|\eta_0\|^{2/a}}
        {\big(\sum_{i=1}^n\|S_i\|^2\big)^{v+n+6+1/2}}.
\]
The two theta weights are $(n+4)L-F$ and $2L-d_{\rm sign}$,
and the auxiliary weight is $a d_{\rm sign}$ for the determinant
sign weight $d_{\rm sign}=(-1^{[p_0]},+1^{[r-p_0]})$.
The effective numerator weight is
\[
 (F+vL)+((n+4)L-F)+(2L-d_{\rm sign})+d_{\rm sign}
               =(v+n+6)L .
\]
At a cut point write \(\varrho\) for the size of the coordinate
tuple \(S_i/w_0\).

If the jets of \(f_0\) through degree \(v+n+6\) vanish, its
squared norm vanishes to order at least \(2(v+n+7)\). The cut rows are
local coordinates, so their sum of squared norms is comparable to
\(\varrho^2\). In this frame,
\[
 Q=O\!\left(\varrho^{\,2(v+n+7)-2(v+n+6+1/2)}\right)
   =O(\varrho)\qquad(\varrho\to0).
\]
At a noncompact end the lower bound for the light sum and the
numerator decay give
\[
  Q\ll 
       \exp\bigl((-4/g_0+2c(v+n+6+1/2)) P\bigr).
\]
For \(v=T_0\) the coefficient of \(P\) is negative by
\(T_0cg_0\le1\) and \(n+6+1/2<T_0\). A positive maximum could only occur
where the factors in the numerator are nonzero and the light sum
is positive.

At such a point the logarithm of the quotient is a
sum of pluriharmonic terms, a negative multiple of
$\log\sum_a|S_a|^2$, and $\varphi_L/2$. The numerator contributes zero Levi form at this point.
The map \(\lambda:z\mapsto[S_1(z):\cdots:S_n(z)]\) has target
dimension \(n-1\) and therefore has a nonzero tangent vector
\(\xi\) in the kernel of its differential. Along \(\xi\) the
pullback of the Fubini--Study form vanishes, so
\[
   \mathcal L(\log\text{quotient})(\xi,\bar\xi)
        =\tfrac12\mathcal L(\varphi_L)(\xi,\bar\xi)>0 .
\]
The last inequality uses the positive curvature of the selected
light weight on the flag chart. A smooth function with a local
maximum has nonpositive Levi form there. Thus all-zero jets force
\(f_0=0\). The dimension at $\kappa$ is
at most $\delta_c\binom{T_0+2n+6}{n}$.

\proofstep{The second upper weight shift.}
The cut contains a point on each component under consideration.
The bounded base theta section of weight $F$ has a nonzero value
at some frame. At that frame the nonvanishing light family gives
\[
  |w_0|^2+\sum_{i=1}^n|S_i|^2>0.
\]
 Multiply the base section by its
$T_0$th power to obtain a nonzero bounded section of weight
$F+T_0L$. The jet injection just proved applies to that bounded
section and would force it to vanish if the cut were empty.

For the upper count at $\kappa+2d$, the weight $(n+4)L-2d$
has positive within-block and cross-block differences. In the
curvature test subtract both $nL$ and the density weight. After both subtractions the remaining cross-block margin satisfies
\[
  7.6T_0-2s_0-2(r-1)>4fr^2.
\]
 Subtracting a further $L$
leaves nonnegative differences. These bounds give bounded
section interpolation on the same cut. Multiplication by one such nonzero
bounded section is injective by holomorphic continuation, and
puts the result at $F+(T_0+n+4)L$.

The maximum argument applies
there with jet degree $T_0+2n+10$. The extra $n+4$ in the
end estimate uses $2n+10.5<T_0$ and $T_0cg_0\le1$, so
$2c(v+n+6.5)<4/g_0$ for $v=T_0+n+4$. Hence the upper count at $\kappa+2d$ is
$\delta_c\binom{T_0+3n+10}{n}$. The component upper bounds
sum by restriction, and the lower bound sums by extension by zero.

For the upper shift we have
\[
 \frac{\binom{T_0+3n+10}{n}}{\binom{T_0+n}{n}}
 \le \exp\left(\frac{n(2n+10)}{T_0+n}\right)<2
 \quad(n\ge2,\ T_0\ge 11n^2).
\]
The loss in the lower shift relative to the unshifted count is
used below. The numbers of cut points and the interpolation
constants disappear from the ratios.
\label{cor:paired-cusp-ranks-9}
\end{proof}


\begin{proposition}[Uniform proportion of the selected character space]
\label{prop:paired-character-share}
In the same two-level model, let \(V[\mu]\) be the character space for \(\mu=(\chi|_U)\psi_0^{-1}\). On the selected family,
\[
 \dim V[\mu]\ge(1-y)\,\frac{R_j}{|U|},\qquad y\le\tfrac18.
\]

Lemma~\ref{lem:paired-character-share} gives the definite case.
For the other lifts use the separating section and the rank bounds
of Lemma~\ref{lem:separating-level-section} and
Proposition~\ref{prop:three-selected-weight-ranks}.
\end{proposition}

\begin{proof}
Write \(P_\mu=\sum_g\mu(g)\rho(g^{-1})\).
Multiplication by a translate of \(u_0\) carries a shifted
cusp section \(X\) into the cusp space of weight \(\kappa\).
This uses boundedness of the translated section.
The map
\[
 X\longmapsto\big(P_\mu(X\cdot h^*u_0)\big)_{h\in U}
\]
from the shifted cusp space to \(|U|\) copies of the
character space is injective. Indeed, at a frame \(z\)
where the translate matrix is invertible, its rows express
the vanishing of those traces as a linear system on the
values \((g^*X)(z)\) times nonzero character and action
scalars. The column for the identity gives \(X(z)=0\).

On every component the determinant is holomorphic and
nonzero somewhere. The invertible locus is dense in its
connected frame domain, so continuity gives vanishing
throughout. Finally,
\[
 \frac{\binom{T_0+n}{n}}{\binom{T_0+2n+6}{n}}
   =\prod_{a=1}^{n}
       \left(1-\frac{n+6}{T_0+n+6+a}\right)
   \ge 1-\frac{n(n+6)}{T_0}.
\]
The family takes \(n(n+6)/T_0\le1/(4(fk)^3)\le1/8\).
This gives the asserted lower share for the character on
both signatures on the indicated two-level model. Taking \(y\) to be the
larger of the two losses gives \(y\le1/8\).
\end{proof}

\paragraph{The rank removed by the arithmetic projector.}

A separating section identifies the levels. The projector can still remove part of their rank, which is counted in the finite representation at the two level primes.

Use the finite group of the two level primes in the same model.
For this calculation twist its action by the inverse of the
central form character $\chi$, so that the chosen character
becomes the product Whittaker character $\lambda$ of $U$. At
each residue factor its value is a nontrivial additive character
of the sum on the first superdiagonal. This twist only multiplies
the scalar in the pullback action. Write $V$ for the cusp space
of weight $\kappa$ with that action and $V_\lambda$ for its
character space. Its dimension is $d_0$.

\begin{lemma}[Trace of the discarded central idempotent]
\label{lem:paired-projector-rank-loss}
For the character space of dimension \(d_0\), write \(r_{\rm bad}\) for the rank discarded by the bad-type projector. The ratios \(\beta_k\) from \eqref{eq:row-bad-share} give
\[
 r_{\rm bad}\le2d_0(\beta_0+\beta_1).
\]
 This comparison uses the integral
arithmetic projector through its scalar extension.
\end{lemma}

\begin{proof}
\proofstep{The finite type decomposition at the two level primes.}
Let $H_k$ be the whole finite linear group at the $k$th level
prime and the upper unitriangulars at the other. There is an
$H_k$-equivariant embedding of $V$ into $a_k$ copies of the
regular representation, with $a_k |H_k|\le 2|U|d_0$.
For a definite form space use the free action on components,
which gives equality with $|U|d_0$ in this bound. In the
indefinite case, let $w$ be a separating section for $H=H_k$.
Write $w_h$ for its ordinary pullbacks, indexed so that the
ring action of $g$ sends $w_h$ to $w_{gh}$. The scalar twist belongs to the action on the form \(v\), while pullback permutes the separating sections. Consequently
\[
  g(w_h v)=w_{gh}\,g(v).
\]

Let $W$ be the invariant subspace of the span of all $w_h v$
for $v\in V$. Consider the map
\[
 V\longrightarrow \operatorname{Map}(H,W),\qquad
 v\longmapsto
    \left(\sum_{g\in H}g(w_h v)\right)_{h\in H}.
\]
For $kv$ the tuple at $h$ becomes the tuple for $v$ at
$k^{-1}h$, so the map intertwines the left regular action.
If the tuple vanishes, the group determinant of the $w_h$
gives $g(v)=0$ for each $g$ on a dense open set of every
component. This gives injectivity. Conversely the map
$(a_h)\mapsto\sum_h w_h a_h$ from $W^{\oplus H}$ has target weight $\kappa+2d$.

If its value vanishes, applying each $g\in H$ gives the same
invertible matrix of translates on the $a_h$ because $W$
consists of invariants. Writing \(a_k=\dim W\), apply first the injection just constructed,
then the three-weight comparison and finally the trace injection for \(U\):
\[
 a_k|H_k|\le\dim V_{\kappa+2d}
 \le2\dim V_{\kappa-d}\le2|U|d_0 .
\]
The bounds on the irreducible multiplicities for $H_k$ are
unchanged by extending the algebraic coefficient field through
a different complex embedding. Indeed, over an algebraic
closure they are the ranks of the central primitive idempotents
divided by the corresponding irreducible degrees. A change of
embedding permutes these idempotents and preserves their degrees.
The regular representation contains each irreducible with
multiplicity its degree. Thus the embedding bound also holds
at the embedding used for the finite type projectors.
\label{cor:paired-good-share-1}

\proofstep{Tracing the central bad projector.}
In the group algebra let $b_k$ be the central idempotent for
bad tuples of semisimple labels at the $k$th prime, extended
by the identity over the other prime. The two commuting central idempotents give the retained projector
\[
  e_{\rm good}=(1-b_0)(1-b_1)=1-b_0-b_1+b_0b_1.
\]
 Let $e_\lambda$
be the normalized character idempotent of $U$.
The operator $b_ke_\lambda$ is idempotent and belongs to
$\mathbb C[H_k]$. The dimension discarded inside $V_\lambda$
is at most the sum of the two ranks of these idempotents.

For the regular representation of $H_k$ the trace of group
algebra multiplication by $B$ is $|H_k|$ times the coefficient
of $1$ in $B$. Indeed each diagonal entry on the group basis
is that coefficient. Consequently, writing $\beta_k$ for that
coefficient in $b_ke_\lambda$ divided by the coefficient in
$e_\lambda$, the rank discarded is at most
\begin{equation}\label{eq:row-bad-share}
 \sum_k a_k\frac{|H_k|}{|U|}\,\beta_k
 \ \le\ 2d_0(\beta_0+\beta_1).
\end{equation}
Here the projector image has its arithmetic meaning even if the
cusp reading uses a different embedding from the type weights.
Both the character kernel and its intersection with the kernel of
the good idempotent are ranks of stacked algebraic linear maps.
This is the use of the two embeddings in the selected comparison.
\label{cor:paired-good-share-2}

\end{proof}

Let \(G=\operatorname{GL}_r(\mathbb F_g)\), \(g=g_k\), with upper
unitriangular subgroup \(U\). For a monic degree-\(r\) polynomial
\(P\) with unit constant coefficient let \(e_P\) be its central
series projector. Write \(e_\psi\) for the idempotent of a standard
nondegenerate character of \(U\) and \([1]\) for the coefficient
of the identity in a group-algebra element.

\begin{lemma}[Local weights of squarefree type labels]
The numbers \(\omega_g(P)=[1](e_Pe_\psi)/[1]e_\psi\) are
nonnegative, sum to one and are independent of the standard
nontrivial additive character used to define \(\psi\). If \(P\)
is squarefree, then
\[
        \omega_g(P)=\prod_{Q\mid P}(g^{\deg Q}-1)^{-1},
\]
where \(Q\) ranges over monic irreducible factors.
\end{lemma}
\begin{proof}
In the regular representation \(\operatorname{tr}_{\rm reg}(a)
=|G|[1]a\); the projectors \(e_Pe_\psi\) split the range of
\(e_\psi\). Conjugation by a diagonal matrix with the prescribed
successive ratios carries one standard character to the other,
fixing each \(e_P\). This proves the first assertions.

For squarefree \(P\) the dual element has distinct eigenvalues;
its centralizer is the torus \(T=\prod_{Q\mid P}\mathbb F_{g^{\deg Q}}^\times\).
The torus character is in general position. The Deligne--Lusztig
degree and Gelfand--Graev multiplicity formulas for \(G\), whose
center is connected, give one irreducible in this series of degree
\(|G|_{g'}/|T|\), with Whittaker multiplicity one. Thus
\[
 \omega_g(P)=\frac{(|G|_{g'}/|T|)\cdot1}{|G|/|U|}=|T|^{-1}.
\]
\cite[\S\S5.24, 7.1, 7.4 and 10.7]{deligne-lusztig-finite}
\end{proof}
At the level prime \(g_k\), sum the products of these local
weights over bad tuples, one factor per split place, to obtain
\(\beta_k\). \label{cor:paired-good-share-3}

\paragraph{Bounds for the four finite label exclusions.}
Lemma~\ref{lem:finite-type-event-bounds} gives the four mass
bounds used here.
The finite type bound used on the selected family is
\begin{align}
 \beta_k\le{}&\frac{3f}{g}+\frac{23 f\log g}{\sqrt g}
     +2fm\left(\frac1{4\min(q,d')}+\frac{1+\log r}{g-1}\right)
       \notag\\
     &+\frac r2(7/8)^f+\exp\left(-\frac{f}{5\log r}\right).
 \label{eq:row-bad-types}
\end{align}
In this formula $q,d'$ are the fixed auxiliary primes at the
point, and $d'$ is distinct from the determinant weight $d$.
The hypotheses are $r\ge100$, $g\ge4096$, distinct odd $q,d'$
with $\max(q,d')\ge16$ and $16qd'\mid g-1$, and the same
character labels as in the projector. The squarefree mass is used
as a subprobability. In degree $r\ge2$ there are $g^r-g^{r-1}$
monic squarefree polynomials by the Euler product, and removing
zero constant terms costs at most $g^{r-1}$.

Since the torus
order of each remaining label is at most $g^r$, the missing
weight is at most $2/g$.
The retained mass is consequently bounded below by
\[
  \frac{g^r-g^{r-1}-g^{r-1}}{g^r}=1-\frac2g .
\]
 A union over the $f$ places fits the
first displayed term. The other terms use four events on the
squarefree labels, namely a pair of roots with equal prime-to-$\ell$
parts at a nonexcluded prime $\ell$, a matching part against a
prescribed scalar or its inverse, a common proper divisor degree
at all places, and the absence of a prime-degree factor of the
specified large size at every place. These are the four tests of
the finite type. We give the counting argument in the next
paragraph.\label{cor:paired-good-share-4}

For the two root tests, let \(g\) be the chosen level prime, let \(r\) be the rank, and put \(N=r\). For a degree-\(j\) label, the order of its multiplicative group and its \(\ell\)-primary part are
\[
  A_j=g^j-1,\qquad A_{j,\ell}=\ell^{v_\ell(g^j-1)}.
\]
 The
squarefree mass of labels that contain specified distinct
irreducibles of degrees $j_1,\ldots,j_u$ is at most
$\prod_i A_{j_i}^{-1}$.

Strip off the specified factors and sum
the remaining weights, which have total at most one in the
remaining degree. This last assertion also follows without the
character formula: if $I_j$ is the number of irreducibles of
degree $j$ with nonzero constant term, then $jI_j\le A_j$ by
counting their distinct roots in $\mathbb F_{g^j}^\times$.
Thus the mass of polynomials with factor counts $m_j$ is at
most $\prod_j(j^{m_j}m_j!)^{-1}$, the probability of that cycle
count under a uniform permutation. This proves domination for
any event specified by the factor degrees.

\begin{lemma}[Mass of the four excluded finite types]
\label{lem:finite-type-event-bounds}
Use the squarefree subprobability measure and the hypotheses of
\eqref{eq:row-bad-types}. Write \(m_{\rm pair},m_{\rm scal},m_{\rm div},m_{\rm cyc}\)
for the probabilities of a repeated prime-to-\(\ell\) root part, a match
with a prescribed scalar or its inverse, a common proper divisor degree,
and absence of the prescribed large prime degree, respectively. Then
\[
 \begin{aligned}
 m_{\rm pair}&\le 23f\log(g)/\sqrt g,&
 m_{\rm scal}&\le2fm\left(\frac1{4\min(q,d')}+\frac{1+\log r}{g-1}\right),\\
 m_{\rm div}&\le\tfrac r2(7/8)^f,&
 m_{\rm cyc}&\le\exp\!\left(-\frac f{5\log r}\right).
 \end{aligned}
\]
The pair test ranges over the nonexcluded \(\ell\), and the last two
bounds use the product over the \(f\) residue factors. The separate
summand \(3f/g\) in \eqref{eq:row-bad-types} allows for the missing
squarefree mass.
\end{lemma}

\begin{proof}
\proofstep{Repeated roots in the squarefree mass.}
When two different roots of a single irreducible of degree $j$
have the same prime-to-$\ell$ part, they are $z,z^{g^a}$ with
$1\le a<j$. Their equality requires $z^{g^a-1}$ to lie in the
$\ell$-primary group. In the cyclic group of order $A_j$, this
allows at most $(g^{\gcd(a,j)}-1)A_{j,\ell}$ choices of $z$.
For a nonexcluded prime, $A_{j,\ell}\le g^{j/4}$ by its
exclusion test.

Since \(\gcd(a,j)\le j/2\), the root count is bounded by
\[
  (g^{\gcd(a,j)}-1)A_{j,\ell}\le g^{j/2}g^{j/4}=g^{3j/4}.
\]
 An equality between distinct roots can occur only
for $\ell\mid A_j$. There are at most $j\log_2 g$ such primes.
Every represented irreducible has at least $j$ winning roots,
because the test is invariant under Frobenius. Dividing the root
count by $j$ and then by $A_j$ bounds the mass for one factor by
\[
 \frac{4096}{4095}\log_2 g
       \sum_{j\ge2}(j-1)u^j,
 \qquad u=g^{-1/4}\le1/8.
\]

For a pair of distinct irreducibles of degrees $i,j$ use the
cyclic groups on the two roots. The common prime-to-$\ell$
part lies in the group of order $g^{\gcd(i,j)}-1$, and the two
primary parts give at most $A_{i,\ell}A_{j,\ell}$ possibilities.
Thus the number of ordered root pairs is at most $g^{3(i+j)/4}$.
The union over the prime factors of $A_i A_j$, followed by
stripping the two factors from the label, bounds the mass by
\[
 \left(\frac{4096}{4095}\right)^2\log_2 g
                \sum_{i,j\ge1}(i+j)u^{i+j}.
\]
The sums are $u^2/(1-u)^2$ and $2u^2/(1-u)^3$. Together they
fit $23\log g/\sqrt g$, with a constant independent of $N$.
This also bounds the whole-place failure by the union over places.
\label{cor:paired-type-mass-1}

\proofstep{A fixed scalar and its root count.}
For a fixed scalar in $\mathbb F_g^\times$, the number of
roots in $\mathbb F_{g^j}^\times$ with matching prime-to-$\ell$
part is at most $A_{j,\ell}$. It suffices to take the union over
primes dividing $A_j$. If a match occurs at another prime the
root is the scalar itself, which is already counted. Divide by
$j$ for the number of roots of each irreducible, and by $A_j$
for the mass of labels containing it.
For fixed \(j,\ell\), these two divisions bound the mass from a fixed scalar by
\[
  \frac{1}{A_j}\,\frac{A_{j,\ell}}{j} .
\]
 At $j=1$, the sum of the
primary parts of $g-1$, divided by $g-1$, is at most
$3/(16\min(q,d'))$. Indeed the three parts for $2,q,d'$ have
size at least $16,q,d'$, and the product of any two is at least
$16\min(q,d')$.

For previous sum \(S\ge2\) and product \(P\), adjoining a primary part
\(c\ge2\) gives
\[
 \frac{S+c}{cP}\le\frac S P,
 \qquad cS-(S+c)=(c-1)(S-1)-1\ge0 .
\] At $j\ge2$, the product of the parts of
$A_j$ supported on $g-1$ is at most $j(g-1)$ by the lifting
formula for powers, also at two since $g\equiv1\pmod4$.
The product of the other parts is at most $A_j/(g-1)$.
In each group the sum of its primary parts is at most the
product. Thus the degree $j$ contribution is at most
$(g-1)/A_j+1/(j(g-1))$. Summing the resulting estimate uses both
\[
  \frac{g-1}{g^j-1}=\frac{1}{1+g+\cdots+g^{j-1}}\le g^{1-j},\qquad \sum_{j=2}^N\frac1j\le\log N.
\]
 This proves the scalar term
in \eqref{eq:row-bad-types} after summing over the $2fm$
scalars and inverse scalars.
\label{cor:paired-type-mass-2}

\proofstep{The divisor-degree and cycle tests.}
A divisor of degree $a$ for a squarefree label corresponds to
an invariant $a$-subset of a permutation in the degree-count
bound. For $0<a<N$, count stable $a$-subsets with weight eight,
except give weight four when the permutation fixes a point on
each side. For a fixed subset \(A\), let \(E_A\) be the event that
it is stable and \(F_A\) the event of a fixed point on
each side. Given \(E_A\), the two restrictions are
independent uniform permutations. The derangement formula
gives probability at least \(1/2\) of a fixed point on
each side, including a side of size one. Thus the mean
contribution from this subset is
\[
 \mathbb P(E_A)(8-4\mathbb P(F_A\mid E_A))\le7\mathbb P(E_A).
\]

When some subset is
stable the total weight is at least eight. In the weight-four
case, swap a fixed point on each side to obtain a different
stable subset of the same size. Summing the stabilization
probability over subsets gives one. Hence the chance of a
stable $a$-subset is at most $7/8$. Independence of the places
and a union over $1\le a\le N/2$ give the fourth term of
\eqref{eq:row-bad-types}.

A permutation can have at most one cycle of prime length in
$(N/2,N-3]$. For each such \(p\), the choice of its support, its cyclic ordering and the permutation on the complement gives
\[
  \Pr(\text{a cycle of length }p)=\binom Np\frac{(p-1)!(N-p)!}{N!}=\frac1p.
\]
 The sum of these reciprocals is at least
$1/(5\log N)$ for $N\ge100$. For $100\le N<16900$ use the sieve to count primes through
$16899$, and sum the lower bounds $10^{-12}\lfloor10^{12}/p\rfloor$.
Multiplying by $5\lfloor\log_2 N\rfloor\cdot0.6931471803$ gives
at least one for every such $N$, using only integer arithmetic.

For larger
$N$, write $n_1=\lfloor N/2\rfloor$, $s_1=\lfloor\sqrt N\rfloor$
and $t=\#\{p:N/2<p\le N-3\}$. The central binomial lower
bound $4^{n_1}/n_1$ and its prime factorization give
\[
 \frac{N-1}{6}\log 4 \le
                 \big(\pi(\sqrt{2n_1})+t+4\big)\log N.
\]
Indeed the primes at most $\sqrt{2n_1}$ cost at most $2n_1$
each, those between $\sqrt{2n_1}$ and $2n_1/3$ cost at most
the primorial up to $2n_1/3$, bounded by $4^{\lfloor2n_1/3\rfloor}$,
and those between $2n_1/3$ and $n_1$ do not divide the binomial.
At most three primes near the upper endpoint were omitted from
the list for $t$.

Now $\pi(\sqrt{2n_1})\le(s_1+9)/3$ by counting
classes coprime to $6$, and $\log N\le(s_1+1)/64+7.70407$ by
$\log u\le u-1$ applied to $u=(s_1+1)/128$. With $s_1\ge130$,
expanding the two sides gives
\[
 (s_1+21)((s_1+1)/64+7.70407)+0.693148\le0.093147s_1^2.
\]
Since $\log 2>0.69314718$, substitution gives $5t\log N\ge N$.
Each reciprocal in the sum is at least $1/N$. The cycle-count domination gives
$1-1/(5\log N)$ as an upper bound for the single-place mass
without one of the prescribed factors. Raising to the $f$th
power bounds the final event in \eqref{eq:row-bad-types}.
\label{cor:paired-type-mass-3}
\end{proof}


\begin{proposition}[A uniform bound on the removed rank]
\label{prop:paired-good-rank}
Use either run of Proposition~\ref{prop:linked-comparison-rows}
and its retained projector in
Proposition~\ref{prop:equivariant-selected-row}. Assume the
finite labels of Lemma~\ref{lem:finite-type-event-bounds} and
size bounds \eqref{eq:paired-bad-type-sizes}. Writing \(d_{\rm rem}\) for the removed rank at the lift, one has
\[
 \delta=\exp(-x^{e_z/7}),\qquad d_{\rm rem}\le\delta d_0,\qquad (1-\delta)d_0\le D_0.
\]
 The finite type bounds are applied
with the same parameters as the central projector trace.
\end{proposition}

\begin{proof}
\proofstep{Summing the finite tests over the selected labels.}

The size bounds for the type mass are
\begin{equation}\label{eq:paired-bad-type-sizes}
 f\ge x^2,\qquad f,m,r\le x^5,\qquad
 g\ge \exp(8x^\epsilon),\qquad
 \min(q,d')\ge\exp(x^{e_z/6}),\qquad
 0<e_z\le\epsilon<1 .
\end{equation}
The exponents \(e_z,\epsilon\) are fixed for the family. On these bounds every term in \eqref{eq:row-bad-types},
including those with $\log g$, is $o(\exp(-x^{e_z/7}))$, uniformly
in the place, lift and selected point. For the logarithm term use
$\log g/\sqrt g\le O(g^{-1/4})$. For the last two terms use
$\log r\le5\log x$ and $f\ge x^2$. At each level prime the resulting family bound is
\[
  4\beta_k\le\delta,\qquad \delta=\exp(-x^{e_z/7})\le\tfrac14.
\]
This threshold is common to the two runs and all the readings.
By \eqref{eq:row-bad-share} the removed rank is at most
$\delta d_0$. Using the tied rank equality from
Section~\ref{subsec:selected-cusp-reading} at each lift gives
$(1-\delta)d_0\le D_0$ on the retained row. Its numerator
comparisons use the corresponding eigenspaces in the same row.
\label{cor:paired-good-share-5}
\end{proof}

